\section{A flat cylinder and a special Lagrangian torus}
\label{sec:completion}

We continue the proof of Theorem~\ref{thm:torus}. All limits are taken along
the subsequence selected in the preceding section. We write $\mu$ for the
limit Riemannian volume and
\[
        H=\C^d/(2\pi i\Z^d),\qquad N=2d.
\]
The standard flat metric on $H$ is again denoted by $\gamma$.
Our goal is to identify $Z$ with this cylinder, obtain smooth convergence
in the actual logarithmic coordinates, and perturb an angular torus.
The possible singular set $S=Z\setminus G$ must be handled first:
flatness on $G$ alone does not give convergence over an entire torus.

\subsection{The coordinate map and its derivative bounds}

On every compact subset of the regular set $G$ the actual logarithmic
chart is defined for all sufficiently small $t$, and $\rho=1$ there.
The functions
\[
       x_j-U_j(p_t),\qquad j=1,\ldots,d,
\]
are harmonic on these subsets. Their local $L^2$ bounds, from
\eqref{eq:centered-poincare}, and interior elliptic estimates give smooth
subsequential limits $u_j$. The angular coordinates have values in the
compact circles $\R/(2\pi\Z)$ and satisfy
$\dd\arg z_j=-\dd x_j\circ I_t$, where $I_t$ is the fiber complex
structure. A diagonal subsequence therefore also gives smooth
circle-valued limits $\vartheta_j$. Thus
\begin{equation}
 F=(u_1+i\vartheta_1,\ldots,u_d+i\vartheta_d):G\longrightarrow H
 \label{eq:limit-map}
\end{equation}
is holomorphic. Diagonal extraction here and below is over a countable
exhaustion of the smooth manifold $G$.

The coordinate density in \eqref{eq:coordinate-density} converges to
a positive constant. Consequently the absolute real Jacobian of $F$,
computed using the limit metric and $\gamma$, is the constant
$J_0=\vartheta^{-1}>0$.
In particular $F$ is a local biholomorphism. It is also injective.
Indeed, if two distinct regular points had the same image, choose
disjoint small regular neighborhoods of these points on which $F$ is
invertible. Smooth convergence and the inverse function theorem imply
that their approximating logarithmic charts both contain one fixed
smaller neighborhood of the common image. This contradicts the
injectivity of the logarithmic chart on each nearby fiber.

For every $R>0$, exhaustion of $B(p,R)\cap G$ by regular compact subsets,
smooth convergence, and the estimates at $p_t$ give, with a fixed
constant $C$,
\begin{align}
 \int_{B(p,R)\cap G}|u_j|^2\,\dd\mu
       &\le CR^2\mu(B(p,2R)),\label{eq:linear-L2}\\
 \int_{B(p,R)\cap G}|\dd u_j|^2\,\dd\mu
       &\le C\mu(B(p,2R)).\label{eq:finite-energy}
\end{align}
Enlarging the radii slightly before taking limits avoids any issue
with boundaries of balls.

To extend these functions across $S$, we use the Sobolev structure
on the whole noncollapsed Ricci limit. Pointed measured
Gromov--Hausdorff convergence implies pointed measured Gromov
convergence by \cite[Proposition~3.30]{GigliMondinoSavare}.
For finite $N>2$, Lott--Villani's \cite[Corollary~E.44]{LottVillani}
preserves nonnegative $N$-Ricci curvature under pointed measured
Gromov--Hausdorff convergence; in particular, the limit satisfies
$\mathrm{CD}(0,N)$. Here $N=2d\ge4$. The finite-dimensional curvature-dimension
condition in the pointed setting is also treated in
\cite[Section~4.2]{GigliMondinoSavare}, while
\cite[Theorem~7.2]{GigliMondinoSavare} preserves quadratic
Cheeger energy. These statements allow measures finite on bounded
sets, without a finite-total-volume assumption. The limit measure
has full support: every open set contains a regular
ball of positive Riemannian volume. Together with
\cite[Theorem~3.17(i)]{ErbarKuwadaSturm}, these stability results
give the $\mathrm{RCD}(0,N)$ structure used below. At zero curvature
the reduced and ordinary curvature-dimension conditions coincide.

On a sufficiently small regular ball, short minimizing $Z$-segments
stay in a larger regular ball. Ambient and Riemannian distances
therefore agree locally. Restrictions of Lipschitz
approximants and classical weak lower semicontinuity give the
lower bound for the relaxed Sobolev energy; smooth compactly
supported approximants, extended by zero, give the upper bound.
Locality and exhaustion therefore identify this energy on $G$
with ordinary Riemannian Sobolev energy.

\begin{lemma}\label{lem:capacity-extension}
Each $u_j$ extends to a weakly harmonic function in
$W^{1,2}_{\mathrm{loc}}(Z)$. Moreover $\lvert\dd u_j\rvert\le C$ on $G$,
with a constant independent of the point. Both $\dd F$ and its local
inverse derivatives are uniformly bounded.
\end{lemma}
\begin{proof}
\emph{Sobolev extension.}
Recall that $\dim_{\mathrm H}S\le N-4$ and
$\mu(B(q,r))\le\omega_Nr^N$. On a fixed bounded subdomain, cover its
compact singular portion by balls of radii $r_i$ with
$\sum_i r_i^{N-2}$ arbitrarily small, and take a finite subcover.
A Lipschitz cutoff on a doubled
covering ball has squared gradient integral at most $C r_i^{N-2}$.
Taking minima of the complementary cutoffs, and exhausting bounded
subdomains, gives cutoffs $\chi_k$ which vanish near $S$, converge
to one off $S$, satisfy $0\le\chi_k\le1$, and have
$\int|\dd\chi_k|^2\to0$. Their exceptional supports can also be
chosen to have measure tending to zero. All these assertions are local;
one inserts a fixed exterior cutoff when necessary.

Let $T_A(u_j)=\max(-A,\min(A,u_j))$. For fixed $A$, the products
$\chi_kT_A(u_j)$ are Sobolev functions across $S$, after extension
by zero there. To justify this before any global path-distance
assertion, insert a fixed compactly supported exterior cutoff.
For each fixed $k$ the product has compact support inside $G$.
The local agreement of distances just noted, applied on a finite
regular cover of the support, shows that the product is locally
ambient Lipschitz;
boundedness controls pairs separated by a fixed positive distance.
It is therefore an ambient Sobolev function, with its ordinary
Riemannian energy on $G$. The extra derivative term has norm at most
$A\|\dd\chi_k\|_2$. Weak compactness and lower semicontinuity give a
Sobolev extension of $T_A(u_j)$ with the energy bound inherited from
$G$. Letting $A\to\infty$ and using
\eqref{eq:linear-L2}--\eqref{eq:finite-energy} gives the claimed
Sobolev extension of $u_j$.
For a compactly supported Lipschitz function $\varphi$, harmonicity
on $G$ gives
\[
 0=\int_G \chi_k\langle\dd u_j,\dd\varphi\rangle\,\dd\mu
       +\int_G\varphi\langle\dd u_j,\dd\chi_k\rangle\,\dd\mu .
\]
The second integral tends to zero by Cauchy--Schwarz, using
\eqref{eq:finite-energy}. The first converges to the integral over
$Z$, since $S$ has zero measure. This proves weak harmonicity.

\emph{Global derivative bounds.}
We use the local harmonic gradient estimate on noncollapsed
Ricci limits, or equivalently its $\mathrm{RCD}(0,N)$ formulation
\cite[Theorem~1.2 and Theorem~4.1(a)]{HuaKellXia}:
\begin{equation}\label{eq:weak-gradient}
 \sup_{B(q,R/2)}|\dd u|
 \le \frac{C(N)}{R}
       \left(\frac{1}{\mu(B(q,2R))}
                    \int_{B(q,2R)}|u|^2\,\dd\mu\right)^{1/2}.
\end{equation}
For this signed-function version, the mean-value estimate applied
to $u^2$ bounds $\sup_{B(q,R)}|u|$ by a constant times the
normalized $L^2$ norm on $B(q,2R)$. Add twice this supremum and
apply the positive-harmonic gradient estimate on the smaller ball;
if the supremum vanishes, the conclusion is immediate. Fix $q$ and
take $R\ge 2d_Z(p,q)$. Then
$B(q,2R)\subset B(p,3R)$ and $B(p,6R)\subset B(q,7R)$,
so doubling compares $\mu(B(q,2R))$ with $\mu(B(p,6R))$.
Equation~\eqref{eq:linear-L2} bounds the right side of
\eqref{eq:weak-gradient} by a constant independent of $q$ and $R$.
Thus $|\dd u_j|\le C$ on $G$.

Holomorphicity bounds the angular derivatives as well, so that
$|\dd F|\le A$ for a fixed $A$. Since the product of its $N$ real
singular values is $J_0>0$, its smallest singular value is at least
$J_0/A^{N-1}$. This gives the inverse derivative bound.
\end{proof}

\subsection{The metric completion is the full cylinder}

We have an injective local biholomorphism on $G$ with uniformly bounded
derivative and inverse derivative.  To extend this identification over
$Z$, we first compare the intrinsic distance on $G$ with the limit
distance, and then continue the inverse along paths in the cylinder.

\begin{lemma}\label{lem:regular-paths}
The intrinsic path distance on $G$ equals the restriction of $d_Z$.
In particular $F$ extends to a Lipschitz map $F:Z\to H$.
\end{lemma}
\begin{proof}
Fix $a,b\in G$ and small balls $A,B$ about them whose closures are
contained in $G$. All minimizing segments between these balls lie
in one compact ball. Cover its singular portion by sufficiently
small balls $B(q_i,r_i)$, disjoint from $A\cup B$ even after doubling,
such that $\sum_i r_i^{N-1}$ is as small as desired. Put
\[
             f=\sum_i r_i^{-1}{\bf1}_{B(q_i,2r_i)}.
\]
The volume upper bound gives
$\int f\,\dd\mu\le C\sum_i r_i^{N-1}$.
The segment inequality for nonnegative integrable functions on Ricci
limits \cite[Theorem~2.15]{CheegerColdingIII} bounds the average, over
endpoints in $A,B$, of the infimum of $\int_\gamma f$ over minimizing
segments $\gamma$ between them. Since $\int f$ can be made arbitrarily
small, some endpoints have this infimum less than one. Choose a
minimizing segment for which $\int_\gamma f<1$; the infimum need not
be attained.

A segment meeting $B(q_i,r_i)$, with its endpoints outside
$B(q_i,2r_i)$, spends length at least $r_i$ in the doubled ball;
it would have $\int f\ge1$. The chosen segment therefore misses
$S$. Short paths inside $A,B$ join it to $a,b$. Letting their radii
tend to zero proves that the path-distance infimum in $G$ is at
most $d_Z(a,b)$; the reverse inequality holds for every path.
The derivative bound of Lemma~\ref{lem:capacity-extension} now
makes $F$ Lipschitz for $d_Z$. The regular set is dense and $H$
is complete, so $F$ extends to $Z$.
\end{proof}

We use two elementary consequences of the small dimension of a removed
set.  First, in a Euclidean ball or flat cylinder, if
$\dim_{\mathrm H}E<N-1$, any two points outside $E$ can be joined
outside $E$ by polygonal paths of length arbitrarily close to their
distance.  The set $E$ need not be closed: intermediate vertices can
be chosen arbitrarily near a shortest segment, outside the cones
of rays through $E$ from the endpoints. Each bad cone is a
countable union of Lipschitz images of $E$ times an interval, and
has dimension at most $\dim_{\mathrm H}E+1<N$. One works in
Euclidean lifts for the cylinder.

Second, if $E$ is closed and $\dim_{\mathrm H}E<N-2$, the complement of
$E$ in a convex Euclidean ball is simply connected. Indeed a loop
has positive distance from $E$ and can be approximated, through a
homotopy in the complement, by a polygonal loop. Cone its finitely
many edges to a generic point of the ball. For each edge, the bad
cone centers lie in countably many bounded Lipschitz images of
$E$ times the edge times a real interval: writing an intersection
point as $(1-t)x+tc$ solves for the center $c$, and one restricts
to $t$ bounded away from zero. This bad set has dimension at most
$\dim_{\mathrm H}E+2<N$. A center outside all the bad sets gives
the required filling.

\begin{lemma}\label{lem:full-cylinder}
The extension $F:Z\to H$ is a bi-Lipschitz homeomorphism.
\end{lemma}
\begin{proof}
Set $E=F(S)$. Its Hausdorff dimension is at most $N-4$; we will prove
closedness only after constructing the inverse on all of $H$.
Local openness of $F$ on $G$ supplies a point of $F(G)\setminus E$.
Start there and lift any finite polygonal path in $H\setminus E$
by the local inverse on $G$. The lift cannot end at a finite
parameter: the inverse derivative bound makes it Cauchy, and
completeness supplies a limit $z\in Z$. Its image is the prescribed
point of the path. This point is outside $F(S)$, hence $z\in G$,
where the local inverse continues the lift. Injectivity on $G$
gives uniqueness. The first avoidance fact shows that
$H\setminus E\subset F(G)$ and that its inverse $K$ obeys
\[
       d_Z(K(a),K(b))\le C\,d_H(a,b)
           \quad(a,b\in H\setminus E).
\]

Both the domain and the image of $K$ are dense. The domain is
dense by the dimension bound. To see density of the image in
$Z$, use local openness of $F$ on each regular neighborhood and
the density of $G$. Completing $K$ thus gives a Lipschitz map
$K:H\to Z$. The compositions $FK$ and $KF$ equal the identity on
dense subsets and hence everywhere. This proves the assertion.
It also makes $E$ closed, since $S$ is closed.
\end{proof}

\subsection{Removing the apparent singularities}

The completion is now the whole cylinder.  We next show that its metric
is smooth across $E=F(S)$, which will prove that $S$ was empty and permit
smooth convergence over every compact cylinder tube.

Transfer the metric on $G$ to $H\setminus E$. It is a flat
K\"ahler metric $g_\infty$, uniformly bounded above and below by
positive multiples of $\gamma$. Consider a small convex complex
coordinate ball $D\subset H$. Its complement of $E$ is simply
connected by the second avoidance fact. A unitary parallel
holomorphic coframe therefore exists there:
\[
     \theta_a=\sum_{j=1}^d A_{aj}(w)\,\dd w_j,\qquad
                g_\infty=\sum_a|\theta_a|^2.
\]
The coefficient matrix $A$ and its inverse are uniformly bounded,
by the two-sided metric bounds.

Each bounded holomorphic coefficient extends holomorphically
across $E$. Here is a direct removal argument. Cover the compact
exceptional set in a smaller ball by balls with
$\sum r_i^{N-1}\to0$, take finite subcovers, and use Lipschitz
cutoffs $\chi_k$ vanishing near $E$.
Then $\|\dd\chi_k\|_{L^1}\to0$. Testing the
Cauchy--Riemann equation for the bounded coefficient against
$\chi_k$ times a smooth test form and passing to the limit proves
the distributional Cauchy--Riemann equation on the whole smaller
ball. The distributional equation implies holomorphicity.
The inverse bound and continuity show that the extended matrix
is nonsingular. The identities expressing that the coframe is
closed, and that the resulting metric is flat and K\"ahler,
extend from the dense complement. Thus $g_\infty$ extends as a
smooth flat K\"ahler metric across $E$. Extensions on overlapping
balls agree by density.

The distance of this extended metric is exactly the completion
distance already obtained. In fact any piecewise smooth path in
a coordinate ball can be approximated in length by paths
avoiding $E$, by polygonal approximation, the avoidance argument,
and uniform continuity of the smooth metric. Subdivision gives
the assertion for arbitrary paths. Conversely every path off
$E$ has its original length. Lemma~\ref{lem:regular-paths} and
completion now identify the two distances.

It follows that every point of $Z$ has Euclidean tangent cones,
so every point is regular in the original Einstein-limit sense.
Thus $S$ is empty. The smooth convergence theorem for Einstein
metrics applies at every point of the cylinder, including those
not initially known to be regular.

The extended metric is translation invariant. To see this, lift
it to $\C^d$, where flatness and simple connectivity give a global
parallel unitary holomorphic coframe. Its coefficients in the
coordinates $w_j$ are bounded entire functions, again by the
uniform metric bounds. Liouville's theorem makes them constant.

The smooth limit must now be expressed in the original logarithmic
coordinates. This will let us perturb a torus while retaining its
angular parametrization. Write
\[
 \mathcal L_t(q)=\bigl(x_j(q)-x_j(p_t)+i\arg z_j(q)\bigr)_{j=1}^d
       \in H,
 \qquad Q_r=\{w\in H:|\operatorname{Re}w|<r\}.
\]
Here $r>0$ is fixed, and $\mathcal L_t$ is defined on the part of $Y_t$
in the crossing chart. The real shifts do not alter its angular
coordinates. Normalize $\eta_t$ by a positive scalar to obtain a
holomorphic volume form $\Omega_t$ with the Calabi--Yau normalization
for $g_t$.

\begin{lemma}[Convergence in the original logarithmic coordinates]
\label{lem:actual-log-convergence}
For every fixed $r>0$, the chart $\mathcal L_t$ contains $Q_r$ in its
image for all sufficiently small $t$ along the selected subsequence.
Its inverse $\Psi_t:Q_r\to Y_t$ satisfies
\[
 \Psi_t^*g_t\longrightarrow g_\infty,\qquad
 \Psi_t^*\omega_t\longrightarrow\omega_\infty,\qquad
 \Psi_t^*\Omega_t\longrightarrow\Omega_\infty
       \quad\hbox{in }C^\infty_{\mathrm{loc}}(Q_r),
\]
where the limits are constant tensors and $\Omega_\infty\ne0$.
The angular coordinates of $\Psi_t$ are the original $\arg z_j$.
\end{lemma}
\begin{proof}
The inclusion $Q_r\subset\mathcal L_t(Y_t\cap\text{chart})$ follows
already from $x(p_t)/s\in B_0\Subset\Sigma$: a bounded real shift leaves
all $d+1$ coordinates $z_0,\ldots,z_d$ exponentially small. The issue
is to identify these coordinate inverses with the smooth convergence
maps.

Take smooth local identifications compatible with the pointed
Gromov--Hausdorff convergence, now available at every point of $Z$.
The embedding-map and harmonic-coordinate arguments of
Section~\ref{sec:metric-estimates} apply on each such neighborhood.
They give smooth local subsequential limits of $\mathcal L_t$.
On its intersection with the old regular set, these limits equal $F$.
Indeed, the old and new identifications send each point to points
whose $g_t$-distance tends to zero; on a slightly larger smooth
neighborhood, the interior harmonic estimates bound the derivatives
of the shifted logarithms and their angular maps. Thus the two
limits coincide there. The old regular set is dense, so every local
limit equals $F$ throughout its neighborhood.

The derivative of $F$ is nonsingular. The inverse function theorem
therefore gives smooth convergence of the local inverse maps on
smaller image neighborhoods. Those inverses agree with $\Psi_t$,
because the actual logarithmic chart is injective. A finite cover
of each compact subset of $Q_r$ now proves the asserted convergence
of the metrics and K\"ahler forms. This argument uses the actual
circle-valued functions $\arg z_j$ at every stage; no change of their
angular marking occurs.

Finally, in these coordinates $\Omega_t$ is a positive normalizing
scalar times $f(z)\,\dd w_1\wedge\cdots\wedge\dd w_d$.
The scalar converges to a positive limit by
\eqref{eq:coordinate-density}. On every fixed tube, $z\to o$
exponentially and $f(z)\to f(o)$ with all derivatives in the $w$
coordinates. This gives the asserted nonzero constant limit of
$\Omega_t$. The metric limit is translation invariant by the
preceding argument, so all three limiting tensors are constant.
\end{proof}

\subsection{A special Lagrangian torus in the nearby fibers}

Fix a tube $Q_{2r}$ and use the charts of
Lemma~\ref{lem:actual-log-convergence} to regard its Calabi--Yau data
as tensors on this one fixed domain. We suppress the pullback symbols.
The polarizing line bundle extends over the total space and is
topologically trivial on a sufficiently small contractible coordinate
neighborhood of the maximal crossing. Thus $\omega_t$ is exact on
$Q_{2r}$. Its constant limit $\omega_\infty$ has zero period on every
angular coordinate two-torus. Since its coefficients are constant,
every angular coefficient vanishes. Thus
\[
             T_0=\{\operatorname{Re}w=0\}\subset Q_{2r}
\]
is Lagrangian for $\omega_\infty$. The constant nonzero form
$\Omega_\infty$ has constant phase on $T_0$, making it special
Lagrangian.

We perturb this torus using the persistence argument for
Calabi--Yau torus tubes in \cite[Section~4, Lemmas~4.3--4.5]{Zhang}.
Its linearization is McLean's special Lagrangian deformation operator
\cite{McLean}. Identify normal vector fields along $T_0$ with one-forms
$\alpha$ by contraction with $\omega_\infty$, and use a fixed normal
graph chart $\iota_\alpha:T_0\to Q_r$. Choose a sufficiently large
integer $k$ and $0<\beta<1$. The source space is
\[
 \mathcal V=\{\alpha\in C^{k+1,\beta}(T^*T_0):
               \alpha\perp\mathcal H^1(T_0)\},
\]
where orthogonality and harmonic forms use the flat metric on $T_0$.
All graphs with $\alpha$ in a fixed sufficiently small ball of
$\mathcal V$ lie in $Q_r$.

The periods $\int_{T_0}\Omega_t$ converge to the nonzero flat period.
Multiply each $\Omega_t$, and its limit, by the phase making this
period positive real. These phases converge. For each small graph set
\[
 \mathcal P_t(\alpha)=
   \bigl(\iota_\alpha^*\omega_t,
          \iota_\alpha^*\operatorname{Im}\Omega_t\bigr).
\]
The first component is exact because $\omega_t$ is exact. The second
has zero integral because graph isotopy preserves
$\int_{T_0}\Omega_t$, whose imaginary part vanishes after the phase
choice. Thus the target is the fixed Banach space $\mathcal W$ consisting of exact
$C^{k,\beta}$ two-forms and exact $C^{k,\beta}$ top forms on $T_0$.
These exact-form spaces are closed.

We have $\mathcal P_\infty(0)=0$. Differentiation of pullbacks of
closed forms gives
\begin{equation}\label{eq:SL-linearization}
 D\mathcal P_\infty(0)\alpha
       =\bigl(\dd\alpha,c\,\dd(*\alpha)\bigr),\qquad c\ne0,
\end{equation}
where the star is the Hodge star of $T_0$. The second identity follows
by contraction in an orthonormal Lagrangian frame. This operator is
an isomorphism $\mathcal V\to\mathcal W$. To see this, rewrite its
second equation as $\delta\alpha=f$, with $f$ of mean zero. For an
exact two-form $\xi$, the Hodge Green operator $\mathcal G$ gives
\[
                \alpha=\delta\mathcal G\xi+\dd\mathcal Gf.
\]
The Hodge identities give $\dd\alpha=\xi$ and $\delta\alpha=f$.
The solution is orthogonal to harmonic forms, and the only solution
of the homogeneous equations in $\mathcal V$ is zero. Elliptic
estimates make the inverse bounded, with one derivative of gain.

Lemma~\ref{lem:actual-log-convergence} gives
$\mathcal P_t\to\mathcal P_\infty$ in $C^1$ on a fixed small ball of
$\mathcal V$, with values in $\mathcal W$. In particular,
$\mathcal P_t(0)\to0$, and their derivatives are uniformly close to
the fixed invertible operator \eqref{eq:SL-linearization} after the
ball is made smaller. The inverse function theorem therefore produces
$\alpha_t\to0$ in $C^{k+1,\beta}$ with $\mathcal P_t(\alpha_t)=0$
for every sufficiently late member of the subsequence. This argument
requires convergence of the ambient tensors, not differentiability
with respect to the degenerating parameter. On each connected
Lagrangian graph, the real restriction of
$\Omega_t$ is nowhere zero. Its sign is therefore constant, and
the positive period makes it positive. Thus the graph is calibrated
by $\operatorname{Re}\Omega_t$ and hence minimal. The
minimal-graph equation is elliptic; smoothness of the ambient data
and elliptic regularity therefore make each graph smooth.

The graphs are embedded, and $\iota_{u\alpha_t}$, $0\le u\le1$,
is an isotopy to $T_0$ inside $Q_r$. The actual chart $\Psi_t$
transports this isotopy to the smooth fiber, ending at an angular
coordinate torus in the prescribed crossing chart. The graph is
Lagrangian and has constant volume-form phase, so it is the torus
required by Theorem~\ref{thm:torus}.

The initial sequence was arbitrary. If the conclusion failed for
arbitrarily small nonzero parameters, choose $t_\nu\to0$ so that no
$Y_{t_\nu}$ contained a torus with the stated properties. The preceding
construction gives such tori on a tail of a subsequence, a contradiction.
This proves the assertion for every sufficiently small nonzero
parameter and completes the proof of Theorem~\ref{thm:torus}.
